\documentclass[UTF8]{ctexart}
\usepackage{amsmath,amssymb,bm}
\usepackage{geometry}
\geometry{a4paper,margin=2.5cm}
\title{贝塞尔曲线公式推演}
\author{}
\date{}

\begin{document}
\maketitle

\section{线性插值}

给定两个控制点 $\bm P_0$ 和 $\bm P_1$，参数 $t\in[0,1]$，线性插值为

\begin{equation}
\bm L(t)=(1-t)\bm P_0+t\bm P_1.
\label{eq:bezier-linear}
\end{equation}

\section{二次贝塞尔曲线}

\begin{align}
\bm Q_0(t)&=(1-t)\bm P_0+t\bm P_1,\\
\bm Q_1(t)&=(1-t)\bm P_1+t\bm P_2.
\end{align}

\begin{equation}
\boxed{
\bm B_2(t)
=
(1-t)^2\bm P_0
+2t(1-t)\bm P_1
+t^2\bm P_2
},
\qquad t\in[0,1].
\label{eq:bezier-quadratic}
\end{equation}

\section{三次贝塞尔曲线}

\begin{equation}
\boxed{
\bm B_3(t)
=
(1-t)^3\bm P_0
+3t(1-t)^2\bm P_1
+3t^2(1-t)\bm P_2
+t^3\bm P_3
}.
\label{eq:bezier-cubic}
\end{equation}

\section{$n$ 次贝塞尔曲线}

De Casteljau 递推写为

\begin{align}
\bm P_i^{(0)}(t)&=\bm P_i,\\
\bm P_i^{(r)}(t)
&=
(1-t)\bm P_i^{(r-1)}(t)
+t\bm P_{i+1}^{(r-1)}(t).
\end{align}

一般形式为

\begin{equation}
\boxed{
\bm B_n(t)
=
\sum_{i=0}^{n}
\binom{n}{i}
(1-t)^{n-i}t^i\bm P_i
},
\qquad t\in[0,1].
\label{eq:bezier-general}
\end{equation}

\subsection{端点与导数}

\begin{align}
\bm B_n(0)&=\bm P_0,\\
\bm B_n(1)&=\bm P_n,\\
\bm B_n'(0)&=n(\bm P_1-\bm P_0),\\
\bm B_n'(1)&=n(\bm P_n-\bm P_{n-1}).
\end{align}

由式~\eqref{eq:bezier-general} 可知，曲线是控制点的 Bernstein 加权组合。

\end{document}
